\input python_2_stepper_style.tex

\newcommand{\sEnv}{E}
\newcommand{\sStore}{\Sigma}
\newcommand{\sRef}[1]{\sCode{\#}#1}
\newcommand{\sBlock}[2]{[\![\,#1,\ #2\,]\!]}
\newcommand{\sStep}[4]{#1 \vdash \langle #2,\ \sStore \rangle \Rightarrow \langle #3,\ #4 \rangle}
\newcommand{\sUn}{\bot}

\section*{Introduction}

The \emph{environment stepper} (``e-stepper'') shows how a Python \S 3 or
Python \S 4 program is evaluated by the \emph{environment model} of SICPy,
Chapter 3. Like the substitution stepper (see the specification of the
Python \S 2 stepper), it rewrites the program itself, one small step at a
time, so that the shape of a process stays visible. In contrast to the
substitution stepper, names are \emph{not} replaced by their values at the time
a binding is made. Instead, every part of the program that is under evaluation
has an \emph{environment} $\sEnv$ in which it looks up names, and assignments
change the frames of the environments. This is what makes assignment, local
state and mutable data structures explainable: the substitution model cannot
express them.

The environments and the data structures that the program refers to are
drawn next to the program as in the CSE machine visualization. The
rewriting rules below and the CSE machine describe the same computations: they
agree on the result of every program, and on the kind of error, if any.

\section*{Configurations}

A step shows a \emph{configuration} $\langle P, \sStore \rangle$ of a
\emph{program term} $P$ and a \emph{store} $\sStore$.

\paragraph{Program term.} The program term $P$ is the remaining program: a
sequence of statements. In it, expressions are rewritten to \emph{values}:
\begin{itemize}
\item primitive values: integers, floats, complex numbers, strings,
  \sCode{True}, \sCode{False} and \sCode{None}, shown as in the substitution
  stepper;
\item \emph{references} $\sRef{k}$ to objects in the store: function objects
  and lists. A reference to a named function object is shown by that name;
  other references are shown as a badge with the number of the object
  in the program and in the drawing;
\item built-in functions, shown by their name.
\end{itemize}
A function body (or lambda body) under evaluation is wrapped as an
\emph{environment block} $\sBlock{\sEnv'}{b}$: the body $b$ is evaluated in
the environment $\sEnv'$. The stepper shows it as the body, surrounded by a
bracket that carries the label of $\sEnv'$ (for example \texttt{E3}) in the
colour of the frame $\sEnv'$ in the drawing.

\paragraph{Store.} The store $\sStore$ is everything that the program can
refer to, besides primitive values: frames, function objects and lists.
We do not describe it as a mathematical structure; it is what the drawing next to
the program shows. The parts of the store \emph{refer to each other}:
\begin{itemize}
\item A \emph{frame} binds names to values, or to $\sUn$ (``unassigned''),
  and refers to its \emph{parent} frame. The \emph{global frame} $G$ has no
  parent. We write $\sEnv(x)$ for the value that $\sEnv$ binds to $x$, and
  $\sEnv = \{\, x_1 \mapsto v_1, \ldots \,\}$ for a frame with these bindings.
\item A \emph{function object} $\mathit{Fn}(n, \mathit{params}, b, \sEnv)$ has
  an optional name $n$, parameters, a body $b$, and refers to the frame
  $\sEnv$ in which the function was created.
\item A \emph{list} $\mathit{List}[v_0, \ldots, v_{n-1}]$ refers to its
  elements. A pair is a list of two elements, as in the CSE machine.
\end{itemize}
A value in a binding or a list that is an object is a \emph{reference}
$\sRef{k}$: it \emph{is} the object, not a copy of it. Two references with the
same number refer to the same object, and a change of the object is seen
through every reference to it. This is the difference to the substitution
model, in which every use of a name gets its own copy of the value.

\emph{The store can contain cycles.} As in the substitution stepper, where a
name that occurs in its own value refers cyclically to that value, the parts
of the store may refer to each other in a circle. The most important cycle is
the one of every function defined in a frame: the frame binds the name of the
function to the function object, and the function object refers to the frame in
which it was created. A list can contain a reference to itself after
\sCode{set\_tail}. The rules never copy or unfold such a cycle; they only
create objects and change what a binding or an element refers to.
We write $\sStore[\sRef{k} := o]$ for the store in which $\sRef{k}$ is
a new object $o$, and $\sStore[\sEnv(x) := v]$ for the store in which the binding
of $x$ in $\sEnv$ refers to $v$ (and similarly for list elements).

The built-in functions live in an implicit frame above $G$. It is not drawn.

\paragraph{Current environment.} The \emph{current environment} $\sEnv$ of a
redex is the environment of the innermost environment block that contains
it, or $G$ if there is none. It is not written in the program, because the
environment of a subexpression is the environment of its block. The rules
below are written as
\[
  \sStep{\sEnv}{e}{e'}{\sStore'}
\]
which reads: ``in environment $\sEnv$, the expression $e$, with store
$\sStore$, reduces to $e'$ and changes the store to $\sStore'$''. A store is
only mentioned in a rule if it is changed or inspected.

\paragraph{Initial configuration.} $G$ is empty, and the program is
evaluated in $G$. $G$ grows: a global name is bound when its first
assignment, function definition or import is executed. This ensures
that a program that later redefines a built-in name, such as
\sCode{print(1); print = 0}, can still use the built-in function before the
assignment. In contrast, the frame of a function call is created with all
names that the function assigns to already present and unassigned (see
\textbf{Function-application} below).

\paragraph{Garbage.} A frame or object is \emph{garbage} if it cannot be
reached from the program term $P$ (through environment blocks and
references) or from $G$. Garbage is not removed, but drawn in grey, so that
the effect of returning from a call is visible: the frame of a call stays
alive if a returned function object refers to it, and becomes garbage
otherwise. The stepper offers to remove the garbage from the drawing
(``Clear dead frames'').

\section*{Reduction}

The reducer $\Rightarrow$ and its transitive closure $\Rightarrow^*$ are used
as in the specification of the Python \S 2 stepper. A \emph{contraction} is a
single step; the order of evaluation is the same: expressions are evaluated
from left to right, inner expressions before the expressions that contain
them. We only give the rules that are different from the substitution
stepper, namely those that involve names, functions and data. The rules for
the primitive operators, for \sCode{and}, \sCode{or}, conditional
expressions, built-in functions on primitive values, and for \sCode{if},
\sCode{pass} and \sCode{breakpoint()} are the same, with the current
environment $\sEnv$ added to every premise.

\subsection*{Names}

\sEntry{Lookup}{A name is replaced by the value that it is bound to. The
frame in which the name is looked up is determined by the rules for scopes
below: the first frame on the path from $\sEnv$ to $G$ that determines the name,
then the built-in frame.
Every name in the program is replaced in a step of its own, from left to right,
and the stepper highlights the binding that was read. (Names of built-in
and library functions, which no frame in the drawing binds, are looked up
as part of the step that uses them.)}
\gentzen{
  \sEnv'\ \mbox{is the frame that determines } x \mbox{ from } \sEnv
  \quad \sEnv'(x) = v \neq \sUn
}{
  \sStep{\sEnv}{x}{v}{\sStore}
}

If $\sEnv'(x) = \sUn$, the reduction is stuck with
\texttt{UnboundLocalError} if $\sEnv'$ is $\sEnv$ itself, and with
\texttt{NameError} (``cannot access free variable'') if $\sEnv'$ is an
enclosing frame, as in the CSE machine. If no frame binds $x$, it is a
\texttt{NameError}.

\subsection*{Functions}

\sEntry{Lambda}{A lambda expression is replaced by a reference to a new
function object that remembers the current environment.}
\gentzen{
  k \mbox{ is fresh}
}{
  \sStep{\sEnv}{\sCode{lambda }x_1\sCode{, }\ldots\sCode{, }x_n\sCode{: }e}{\sRef{k}}
  {\sStore[\sRef{k} := \mathit{Fn}(\varepsilon, x_1\ldots x_n, e, \sEnv)]}
}

\sEntry{Function-definition}{A function definition is removed from the
program. A new function object that remembers the current environment is
created, and its name is bound to it in the frame that the name belongs to.}
\gentzen{
  k \mbox{ is fresh}\quad \sEnv' \mbox{ is the frame that } f \mbox{ belongs to}
}{
  \begin{array}{l}
  \sEnv \vdash \langle \sCode{def }f\sCode{(}x_1\sCode{, }\ldots\sCode{, }x_n\sCode{):}\ b\ \sVar{statement}\ldots,\ \sStore \rangle \Rightarrow \\
  \qquad \langle \sVar{statement}\ldots,\ \sStore[\sRef{k} := \mathit{Fn}(f, x_1\ldots x_n, b, \sEnv)][\sEnv'(f) := \sRef{k}] \rangle
  \end{array}
}

\sEntry{Function-application}{The application of a reference to a function
object, to as many arguments as it has parameters, all of which are values,
is replaced by the body of the function in an environment block. The block's
environment is a new frame whose parent is the environment \emph{in which the
function was created} (and not the environment of the caller). The
frame binds the parameters to the arguments (a rest parameter
\sCode{*r} to a new list of the remaining arguments), and every other name
that the body assigns to as unassigned.
Names that the body declares \sCode{global} or \sCode{nonlocal} are not bound
in the new frame. If the number of arguments does not match, the reduction is
stuck with a \texttt{TypeError}, as in the CSE machine.}
\gentzen{
  \sStore(\sRef{k}) = \mathit{Fn}(n, x_1\ldots x_m, b, \sEnv')\quad
  \sEnv'' \mbox{ is fresh with parent } \sEnv'
}{
  \sStep{\sEnv}{\sRef{k}\sCode{(}v_1\sCode{, }\ldots\sCode{, }v_m\sCode{)}}
    {\sBlock{\sEnv''}{b}}
    {\sStore[\sEnv'' := \{x_i \mapsto v_i\} \cup \{y \mapsto \sUn \mid y \in \mathit{assigned}(b)\}]}
}

\sEntry{Block-return}{An environment block whose first statement is a
\sCode{return} reduces to the returned expression, evaluated in $\sEnv$:
the block disappears, the statements after the \sCode{return} are dropped,
and the frame of the block stays in the store (it may be garbage).
A \sCode{return} without an expression returns \sCode{None}.}
\gentzen{}{
  \sBlock{\sEnv''}{\sCode{return }e\ \sVar{statement}\ldots} \Rightarrow e
}

\sEntry{Block-value}{An environment block whose body is a lambda body
that has been reduced to a value $v$ reduces to $v$. An environment block whose
statements are all finished reduces to \sCode{None}.}
\gentzen{}{
  \sBlock{\sEnv''}{v} \Rightarrow v
  \qquad\qquad
  \sBlock{\sEnv''}{\varepsilon} \Rightarrow \sCode{None}
}

\sEntry{Block-reduce}{The first statement of an environment block is
reduced in the \emph{frame of the block}, not in $\sEnv$.}
\gentzen{
  \sStep{\sEnv''}{s}{s'}{\sStore'}
}{
  \sStep{\sEnv}{\sBlock{\sEnv''}{s\ \sVar{statement}\ldots}}
    {\sBlock{\sEnv''}{s'\ \sVar{statement}\ldots}}{\sStore'}
}

A sequence of statements in an environment block is reduced statement by
statement as a program, with environment $\sEnv''$. A consumed statement
disappears. Calls of library functions (such as \sCode{map}) are applied in
a single step.

\subsection*{Assignment}

\sEntry{Assignment}{The right-hand expression of an assignment is evaluated.
When it is a value, the statement is removed from the program, and the name is
rebound to the value in the frame that the name belongs to (see Scopes).
If the frame was created before, this is a change of the store that the
substitution model cannot express.}
\gentzen{
  \sEnv' \mbox{ is the frame that } x \mbox{ belongs to}
}{
  \sStep{\sEnv}{x\sCode{ = }v\ \sVar{statement}\ldots}{\sVar{statement}\ldots}
    {\sStore[\sEnv'(x) := v]}
}

\sEntry{Global-and-nonlocal}{The declarations \sCode{global} and
\sCode{nonlocal} are not statements that run: they have no step of their own.
They remain in the program as source and have been taken into account when
the frame of the function was created.}

\subsection*{Scopes}

The frame that a name $x$ \emph{belongs to}, for both lookup and assignment, is
decided statically by Python's rule, as in the CSE machine:
\begin{itemize}
\item if the function declares \sCode{global x}, it is $G$;
\item if the function declares \sCode{nonlocal x}, it is the nearest
  enclosing function frame that binds $x$;
\item otherwise, if $x$ is a parameter of the function or is assigned to
  anywhere in the body, it is the frame of the function call;
\item otherwise, the name is looked up in the parent frames, from the
  innermost to $G$.
\end{itemize}

\subsection*{Lists and pairs}

\sEntry{List-literal}{A list expression whose elements are all values is
replaced by a reference to a new list.}
\gentzen{
  k \mbox{ is fresh}
}{
  \sStep{\sEnv}{\sCode{[}v_0\sCode{, }\ldots\sCode{, }v_{n-1}\sCode{]}}{\sRef{k}}
    {\sStore[\sRef{k} := \mathit{List}[v_0,\ldots,v_{n-1}]]}
}

\sEntry{List-access}{The access of an element of a list is replaced by the
element. A negative index counts from the end. An index that is out of range
causes an \texttt{IndexError}, and an index that is not an integer, a
\texttt{TypeError}; the reduction is then stuck. Strings can also be accessed;
the result is a string of one character.}
\gentzen{
  \sStore(\sRef{k}) = \mathit{List}[v_0,\ldots,v_{n-1}]\quad
  -n \leq i < n
}{
  \sRef{k}\sCode{[}i\sCode{]} \Rightarrow v_{i \bmod n}
}

\sEntry{List-assignment}{The assignment to an element of a list is removed
from the program, and the element is changed in the store.}
\gentzen{
  -n \leq i < n
}{
  \sStep{\sEnv}{\sRef{k}\sCode{[}i\sCode{] = }v\ \sVar{statement}\ldots}{\sVar{statement}\ldots}
    {\sStore[\sRef{k}[i \bmod n] := v]}
}

\sEntry{Pair-construction}{The application of \sCode{pair}, \sCode{llist}
(and, in Python \S 3, \sCode{stream}) to values is replaced by a reference to
the outermost new list; all the cells are created in the same step.}
\gentzen{
  k_1, \ldots, k_n \mbox{ are fresh}
}{
  \sStep{\sEnv}{\sCode{pair(}v_1\sCode{, }v_2\sCode{)}}{\sRef{k_1}}
    {\sStore[\sRef{k_1} := \mathit{List}[v_1, v_2]]}
}

\sEntry{Pair-mutation}{The application of \sCode{set\_head} or
\sCode{set\_tail} to a reference to a pair and a value is replaced by
\sCode{None}, and the first or second element of the pair is changed in the
store.}
\gentzen{
  \sStore(\sRef{k}) = \mathit{List}[v_0, v_1]
}{
  \sStep{\sEnv}{\sCode{set\_head(}\sRef{k}\sCode{, }v\sCode{)}}{\sCode{None}}
    {\sStore[\sRef{k}[0] := v]}
}

The other built-in functions of the linked list and stream libraries, and the
functions that read lists such as \sCode{head}, \sCode{tail} and
\sCode{is\_pair}, work as in the substitution stepper, on the lists in the store
instead of on list literals.

\sEntry{Identity}{\sCode{v is w} is replaced by \sCode{True} if $v$ and $w$
are the same primitive value or the same reference, and by \sCode{False}
otherwise; \sCode{is not} is the negation.}
\gentzen{}{
  \sRef{k}\sCode{ is }\sRef{k} \Rightarrow \sCode{True}
}

\subsection*{Loops}

\sEntry{While-unfold}{A \sCode{while} loop is unfolded by the textbook rule
before its test is evaluated. The loop stays in the program, in full,
while the test is evaluated; so the program and the store always contain
what is needed to continue by hand. The \sCode{if} is then reduced by the
rules for conditional statements.}
\gentzen{}{
  \sCode{while }e\sCode{:}\ b\ \Rightarrow\ 
  \sCode{if }e\sCode{:}\ b\ \ \sCode{while }e\sCode{:}\ b
}

\sEntry{For-range}{A \sCode{for} loop over \sCode{range(}$a$\sCode{, }$b$\sCode{,
}$s$\sCode{)} first evaluates the arguments of \sCode{range} to values. The
loop is then tested step by step: while the current value $c$ is in the
range, the target is assigned $c$ and the body is run, followed by the loop
with the next value $c + s$; otherwise the loop ends.}
\gentzen{
  c \mbox{ is in the range } a, b, s
}{
  \sCode{for }x\sCode{ in range(}c\sCode{, }b\sCode{, }s\sCode{):}\ \sVar{body}
  \ \Rightarrow\ 
  x\sCode{ = }c\ \ \sVar{body}\ \ \sCode{for }x\sCode{ in range(}c+s\sCode{, }b\sCode{, }s\sCode{):}\ \sVar{body}
}

\sEntry{Break-and-continue}{Within an iteration, \sCode{break} removes the rest of
the iteration and the loop, and \sCode{continue} removes the rest of the
iteration and continues with the loop.}

\section*{Worked example}

\begin{lstlisting}
def make_withdraw(balance):
    def withdraw(amount):
        nonlocal balance
        if balance >= amount:
            balance = balance - amount
            return balance
        else:
            return 'Insufficient funds'
    return withdraw

W1 = make_withdraw(100)
W1(50)
\end{lstlisting}

Initially, $G$ is empty.
\begin{enumerate}
\item \textbf{Function-definition}: the definition of \sCode{make\_withdraw}
  is consumed. A function object $\sRef{1}$ is created
  for \sCode{make\_withdraw}, with $G$ as its frame, and $G$ binds
  \sCode{make\_withdraw} to $\sRef{1}$.
\item \textbf{Lookup}: in \sCode{W1 = make\_withdraw(100)}, the name
  \sCode{make\_withdraw} is replaced by $\sRef{1}$, in a step of its own.
\item \textbf{Function-application}: a new frame
  $E_1 = \{\mathtt{balance} \mapsto 100, \mathtt{withdraw} \mapsto \sUn\}$ with
  parent $G$. The program becomes
  \sCode{W1 = }$\sBlock{E_1}{\sCode{def withdraw}\ldots\sCode{; return withdraw}}$.
\item \textbf{Function-definition} in $E_1$: a function object $\sRef{2}$ is
  created for \sCode{withdraw}, with $E_1$ as its frame, and $E_1$ binds
  \sCode{withdraw} to $\sRef{2}$. This is a cycle: $E_1$ refers to $\sRef{2}$,
  and $\sRef{2}$ refers to $E_1$.
\item \textbf{Lookup}, then \textbf{Block-return}: the block is replaced by
  the value of \sCode{withdraw}, which is $\sRef{2}$ (shown as
  \texttt{withdraw}).
\item \textbf{Assignment}: $G(\mathtt{W1}) = \sRef{2}$. $E_1$ stays alive, since
  $\sRef{2}$ refers to it.
\item \textbf{Function-application}: \sCode{W1(50)} is looked up and applied. A
  new frame $E_2 = \{\mathtt{amount} \mapsto 50\}$ with parent $E_1$ (the frame
  in which $\sRef{2}$ was created, and not the caller's). $E_2$ has no binding
  for \sCode{balance}: \sCode{nonlocal balance} decided that.
\item \textbf{Lookup}: in \sCode{balance >= amount}, \sCode{balance} is found
  in $E_1$ (100), then \sCode{amount} in $E_2$ (50), each in its own step.
  \sCode{100 >= 50} is \sCode{True}, and the \sCode{if} runs its first branch.
\item \textbf{Assignment}: \sCode{balance = balance - amount} reduces to
  \sCode{balance = 50}. With \sCode{nonlocal}, the binding is changed
  \emph{in $E_1$}: this is the step that the substitution model
  cannot express.
\item \textbf{Block-return}: \sCode{return balance} gives 50 and the block
  disappears. $E_2$ is now garbage; $E_1$ is not.
\end{enumerate}
A second call \sCode{W1(50)} creates $E_3$ with parent $E_1$ and finds
\sCode{balance} $= 50$ there.

\section*{Steps and their limit}

Steps are numbered, explained and limited as in the Python \S 2 stepper:
step 0 shows the program before evaluation, each contraction contributes two
steps (the redex highlighted, then the result highlighted), and the last step
is ``Evaluation complete'', ``Evaluation stuck'' or ``Maximum number of steps
exceeded''. The number of the last step does not exceed the step limit.
Each step also shows the store as it is at that step: the frames and objects
of the drawing, with the frame of the current redex highlighted.

A gutter breakpoint in the editor, and the statement \sCode{breakpoint()},
mark the step in which the statement is reached (for a gutter breakpoint, the
first step of the statement); the \emph{previous breakpoint} and \emph{next
breakpoint} buttons jump between these steps.

\section*{Python \S 4}

The e-stepper does not yet support the metacircular evaluator support of
Python \S 4 (\sCode{parse}, \sCode{tokenize} and
\sCode{apply\_in\_underlying\_python}), nor constructs outside the Python \S 3
syntax. When the reduction reaches such a construct, it stops with the
message that the construct is not supported.
